
\documentclass{article}
\usepackage{colortbl}
\usepackage{makecell}
\usepackage{multirow}
\usepackage{supertabular}

\begin{document}

\newcounter{utterance}

\centering \large Interaction Transcript for game `referencegame', experiment `number\_grids', episode 14 with reflexion{-}Qwen3.5{-}9B.
\vspace{24pt}

{ \footnotesize  \setcounter{utterance}{1}
\setlength{\tabcolsep}{0pt}
\begin{supertabular}{c@{$\;$}|p{.15\linewidth}@{}p{.15\linewidth}p{.15\linewidth}p{.15\linewidth}p{.15\linewidth}p{.15\linewidth}}
    \# & $\;$A & \multicolumn{4}{c}{Game Master} & $\;\:$B\\
    \hline

    \theutterance \stepcounter{utterance}  
    & & \multicolumn{4}{p{0.6\linewidth}}{
        \cellcolor[rgb]{0.9,0.9,0.9}{
            \makecell[{{p{\linewidth}}}]{
                \texttt{\tiny{[P1$\langle$GM]}}
                \texttt{You have attempted this same task before and did not succeed. These are the plans you wrote after those attempts.} \\
\texttt{Trial \#1: New plan: My previous description was too vague and failed to explicitly distinguish the target from the distractors by pinpointing the exact row differences. Specifically, I mentioned a "complete rectangular border" which is true for the target, but I did not clearly state that the distractors have a "broken" bottom border (missing the bottom{-}left cell) or a "broken" middle section in the third row. Next time, I will construct the expression by first confirming the shared features (the vertical line in column 2 and the top/bottom rows in the target) and then immediately contrasting them with the specific missing cells in the distractors: explicitly stating that the target has filled cells at (row 3, col 1) and (row 4, col 1) which are empty in both distractors, and that the target has a solid bottom row while Distractor 1 has a gap at the bottom{-}left and Distractor 2 has a gap at the bottom{-}left and a different middle pattern. I will ensure the phrase "unlike the distractors" is followed by the precise coordinates or relative positions of the differing cells.} \\
\\ 
\\ 
\texttt{You are given three grids, where each of them is 5 by 5 in size.} \\
\texttt{Grids have empty cells marked with "▢" and filled cells marked with "X".} \\
\texttt{Your task is to generate a referring expression that best describes the target grid while distinguishing it from the two other distractor grids.} \\
\texttt{The first grid is the target grid, and the following two grids are the distractors.} \\
\\ 
\texttt{Target grid:} \\
\\ 
\texttt{□ X X X X} \\
\texttt{□ X □ □ X} \\
\texttt{□ X □ □ X} \\
\texttt{□ X □ □ X} \\
\texttt{□ X X X X} \\
\\ 
\texttt{Distractor grid 1:} \\
\\ 
\texttt{□ X X X X} \\
\texttt{□ X □ □ X} \\
\texttt{□ X X X X} \\
\texttt{□ X □ □ X} \\
\texttt{□ X X X X} \\
\\ 
\texttt{Distractor grid 2:} \\
\\ 
\texttt{□ X X X X} \\
\texttt{□ X □ □ X} \\
\texttt{□ X X X X} \\
\texttt{□ □ □ □ □} \\
\texttt{□ X X X X} \\
\\ 
\texttt{Instruction: Describe the target grid.} \\
\texttt{Generate the referring expression starting with the tag "Expression: " for the given target grid. Omit any other text.} \\
            }
        }
    }
    & \\ \\

    \theutterance \stepcounter{utterance}  
    & \multicolumn{4}{p{0.6\linewidth}}{
        \cellcolor[rgb]{0.9,0.9,0.9}{
            \makecell[{{p{\linewidth}}}]{
                \texttt{\tiny{[P1$\rangle$GM]}}
                \texttt{Expression: The grid features a vertical line of filled cells in the second column spanning all five rows, with the top row and bottom row completely filled except for the second column's vertical line (forming a rectangular border), unlike the distractors where the third row has a filled cell in the third column (Distractor 1) or the fourth row is entirely empty (Distractor 2), and specifically, the target has filled cells at (row 3, col 1) and (row 4, col 1) which are empty in both distractors, creating a solid bottom{-}left corner that the others lack.} \\
            }
        }
    }
    & & \\ \\

    \theutterance \stepcounter{utterance}  
    & & & \multicolumn{2}{p{0.3\linewidth}}{
        \cellcolor[rgb]{0.9,0.9,0.9}{
            \makecell[{{p{\linewidth}}}]{
                \texttt{\tiny{[GM$|$GM]}}
                \texttt{Expression: The grid features a vertical line of filled cells in the second column spanning all five rows, with the top row and bottom row completely filled except for the second column's vertical line (forming a rectangular border), unlike the distractors where the third row has a filled cell in the third column (Distractor 1) or the fourth row is entirely empty (Distractor 2), and specifically, the target has filled cells at (row 3, col 1) and (row 4, col 1) which are empty in both distractors, creating a solid bottom{-}left corner that the others lack.} \\
            }
        }
    }
    & & \\ \\

    \theutterance \stepcounter{utterance}  
    & & \multicolumn{4}{p{0.6\linewidth}}{
        \cellcolor[rgb]{0.9,0.9,0.9}{
            \makecell[{{p{\linewidth}}}]{
                \texttt{\tiny{[GM$\rangle$P2]}}
                \texttt{You have attempted this same task before and did not succeed. These are the plans you wrote after those attempts.} \\
\texttt{Trial \#1: New plan: I incorrectly identified the "Second" grid because I focused on the vertical line in the second column but failed to verify the horizontal border requirement. The expression demands a complete rectangular border (filled first and last rows) with an empty interior except for the second column. The "Second" grid has empty cells in the first and last rows (specifically columns 3 and 4), violating the border condition. The "Third" grid, however, has filled cells in the entire first and last rows, creating the required complete border, while maintaining the vertical line in the second column and an empty interior elsewhere. Next time, I will strictly check for the presence of filled cells in all four corners and the entirety of the first and last rows before confirming a match.} \\
\texttt{Answer: Third} \\
\\ 
\\ 
\texttt{You are given three grids, where each of them is 5 by 5 in size.} \\
\texttt{Grids have empty cells marked with "▢" and filled cells marked with "X".} \\
\texttt{You are also given a referring expression that describes one of the given grids.} \\
\texttt{Your task is to select the grid that matches the given referring expression.} \\
\texttt{Generate only the number (in text) of the grid that matches the given expression by selecting first, second, or third.} \\
\\ 
\texttt{First:} \\
\\ 
\texttt{□ X X X X} \\
\texttt{□ X □ □ X} \\
\texttt{□ X X X X} \\
\texttt{□ □ □ □ □} \\
\texttt{□ X X X X} \\
\\ 
\texttt{Second:} \\
\\ 
\texttt{□ X X X X} \\
\texttt{□ X □ □ X} \\
\texttt{□ X X X X} \\
\texttt{□ X □ □ X} \\
\texttt{□ X X X X} \\
\\ 
\texttt{Third:} \\
\\ 
\texttt{□ X X X X} \\
\texttt{□ X □ □ X} \\
\texttt{□ X □ □ X} \\
\texttt{□ X □ □ X} \\
\texttt{□ X X X X} \\
\\ 
\texttt{Expression: The grid features a vertical line of filled cells in the second column spanning all five rows, with the top row and bottom row completely filled except for the second column's vertical line (forming a rectangular border), unlike the distractors where the third row has a filled cell in the third column (Distractor 1) or the fourth row is entirely empty (Distractor 2), and specifically, the target has filled cells at (row 3, col 1) and (row 4, col 1) which are empty in both distractors, creating a solid bottom{-}left corner that the others lack.} \\
\texttt{Question: Which grid does the expression refer to?} \\
\texttt{Start with the tag "Answer: ", followed by your selection. Omit any other text.} \\
            }
        }
    }
    & \\ \\

    \theutterance \stepcounter{utterance}  
    & & & \multicolumn{4}{p{0.6\linewidth}}{
        \cellcolor[rgb]{0.9,0.9,0.9}{
            \makecell[{{p{\linewidth}}}]{
                \texttt{\tiny{[GM$\langle$P2]}}
                \texttt{Answer: First} \\
            }
        }
    }
     \\ \\

    \theutterance \stepcounter{utterance}  
    & & & \multicolumn{2}{p{0.3\linewidth}}{
        \cellcolor[rgb]{0.9,0.9,0.9}{
            \makecell[{{p{\linewidth}}}]{
                \texttt{\tiny{[GM$|$GM]}}
                \texttt{first} \\
            }
        }
    }
    & & \\ \\

    \theutterance \stepcounter{utterance}  
    & \multicolumn{4}{p{0.6\linewidth}}{
        \cellcolor[rgb]{0.95,0.95,0.95}{
            \makecell[{{p{\linewidth}}}]{
                \texttt{\tiny{[P2$\rightarrow$P1]}}
                \texttt{Post{-}episode Reflexion (not sent to GM):} \\
\texttt{New plan: My previous description failed because it relied on vague terms like "solid bottom{-}left corner" and "broken border" without explicitly mapping the specific cell coordinates that differ between the grids. The target grid is unique because it has filled cells at (row 3, col 1) and (row 4, col 1), whereas Distractor 1 has a filled cell at (row 3, col 3) and Distractor 2 has an entirely empty fourth row. Next time, I will construct the expression by first stating the shared vertical line in column 2, then immediately contrasting the target's specific filled cells at (3,1) and (4,1) against the specific differing features of the distractors: explicitly noting that Distractor 1 has a filled cell at (3,3) instead of (3,1), and Distractor 2 has an empty row 4 instead of the filled cells at (4,1). I will ensure every distinguishing feature is tied to exact row and column coordinates to eliminate ambiguity.} \\
            }
        }
    }
    & & \\ \\

    \theutterance \stepcounter{utterance}  
    & & & \multicolumn{4}{p{0.6\linewidth}}{
        \cellcolor[rgb]{0.95,0.95,0.95}{
            \makecell[{{p{\linewidth}}}]{
                \texttt{\tiny{[P1$\rightarrow$P2]}}
                \texttt{Post{-}episode Reflexion (not sent to GM):} \\
\texttt{New plan: I incorrectly selected "First" because I misread the grid layout and failed to verify the specific row conditions mentioned in the expression. The expression explicitly states that the target has filled cells at (row 3, col 1) and (row 4, col 1), which are empty in the "First" grid (where row 3 and 4 are entirely empty). The "Second" grid matches this description perfectly: it has a vertical line in column 2, empty rows 3 and 4 in column 1, and a filled bottom row. Next time, I will cross{-}reference every specific coordinate mentioned in the expression (like row 3, col 1) against the actual grid content before making a selection, ensuring I don't overlook the "empty interior" constraint that distinguishes the correct grid from the distractors.} \\
\texttt{Answer: Second} \\
            }
        }
    }
     \\ \\

\end{supertabular}
}

\end{document}
